% GENERATED FILE. Edit canonical lessons, then run npm run generate:books.
% canonical-source-hash: fnv1a64:8fad949ce3a40640
% canonical-lessons: BN-C407-baja, BN-C407-beyara, BN-C407-menu, BN-C407-bokshish, BN-C407-londri, BN-R407-review-400-407

\chapter{Bells, Menus and Tips}
\label{ch:bn-bells-menus-and-tips}
\hlchaptermodality{407}

\begin{chapteropening}
\textbf{By the end of this chapter:} \emph{I can hear the bell, read the menu, tip the bearer and send clothes to the laundry.}

Complete the last lesson of chapter 407: I can hear the bell, read the menu, tip the bearer and send clothes to the laundry.
\end{chapteropening}

\section[\texorpdfstring{\bn{বাজা} (bājā)}{বাজা (bājā)}]{\bn{বাজা} (bājā) — to ring; to sound; to strike (the hour)}

\label{lesson:BN-C407-baja}



\begin{quote}
Before the new one: say the Bengali for a place to stay, then the Bengali for to tap (on a door).
\end{quote}

\subsection*{You'll want to know: \bn{বাজা}}
\textbf{\bn{বাজা}} — \emph{bājā} — \textquotedblleft{}to ring; to sound; to strike (the hour)\textquotedblright{}.

\begin{quote}
\bn{স্টেশনে} \bn{ঘণ্টা} \bn{বাজল}, \bn{আর} \bn{ট্রেন} \bn{ছাড়ল।}
\end{quote}

\begin{quote}
\emph{sṭeshône ghônṭā bājlô, ār ṭren chhāṛlô.}
\end{quote}

\begin{quote}
The bell rang at the station, and the train left.
\end{quote}

Past (he/she …): \bn{সেই} \bn{ঘণ্টা} \bn{বাজল}, \emph{sei ghônṭā bājlô}.

Also of the hour: \bn{দশটা} \bn{বাজল}, it struck ten. \bn{বাজানো}, to play an instrument, is the causing partner.

\begin{cousinweb}
\textbf{\bn{বাজা}} (\emph{bājā}) is inherited: CDIAL lists Bengali \emph{bājā}, \textquotedblleft{}to sound\textquotedblright{}, of a musical instrument, under Sanskrit \emph{vādyatē}, \textquotedblleft{}is sounded, is played\textquotedblright{} (CDIAL 11513).
\end{cousinweb}

\begin{grammarlens}[title={What you've built}]
One more word for getting around.
\end{grammarlens}

\subsection*{Your turn}
\begin{itemize}
\raggedright
  \item \emph{Say it:} \emph{bājā}
  \item \emph{Say it:} \emph{bājā}, once more
  \item \emph{Say it:} \emph{sṭeshône ghônṭā bājlô, ār ṭren chhāṛlô.} — The bell rang at the station, and the train left.
  \item \emph{Recall:} say the Bengali for a place to stay, then the Bengali for to tap (on a door), then say \emph{bājā} again
\end{itemize}

\begin{tcolorbox}[breakable,colback=teal!4,colframe=teal!35!black,title={Before you move on}]
What does \bn{বাজা} mean? (\textquotedblleft{}To ring; to sound; to strike (the hour)\textquotedblright{}.) Say it once more.
\end{tcolorbox}
\section[\texorpdfstring{\bn{বেয়ারা} (beyārā)}{বেয়ারা (beyārā)}]{\bn{বেয়ারা} (beyārā) — a bearer; a waiter (in a hotel)}

\label{lesson:BN-C407-beyara}



\begin{quote}
Before the new one: say the Bengali for to tap (on a door), then the Bengali for to ring.
\end{quote}

\subsection*{You'll want to know: \bn{বেয়ারা}}
\textbf{\bn{বেয়ারা}} — \emph{beyārā} — \textquotedblleft{}a bearer; a waiter (in a hotel)\textquotedblright{}.

\begin{quote}
\bn{বেয়ারা} \bn{আমাদের} \bn{ঘরে} \bn{চা} \bn{নিয়ে} \bn{এল।}
\end{quote}

\begin{quote}
\emph{beyārā āmāder ghôre chā niye elô.}
\end{quote}

\begin{quote}
The bearer brought tea to our room.
\end{quote}

\begin{cousinweb}
\textbf{\bn{বেয়ারা}} (\emph{beyārā}) is usually traced to English \emph{bearer}, the old word in Indian hotels and clubs for a servant who carries and serves.
\end{cousinweb}

\begin{grammarlens}[title={What you've built}]
One more word for getting around.
\end{grammarlens}

\subsection*{Your turn}
\begin{itemize}
\raggedright
  \item \emph{Say it:} \emph{beyārā}
  \item \emph{Say it:} \emph{beyārā}, once more
  \item \emph{Say it:} \emph{beyārā āmāder ghôre chā niye elô.} — The bearer brought tea to our room.
  \item \emph{Recall:} say the Bengali for to tap (on a door), then the Bengali for to ring, then say \emph{beyārā} again
\end{itemize}

\begin{tcolorbox}[breakable,colback=teal!4,colframe=teal!35!black,title={Before you move on}]
What does \bn{বেয়ারা} mean? (\textquotedblleft{}A bearer; a waiter (in a hotel)\textquotedblright{}.) Say it once more.
\end{tcolorbox}
\section[\texorpdfstring{\bn{মেনু} (menu)}{মেনু (menu)}]{\bn{মেনু} (menu) — a menu}

\label{lesson:BN-C407-menu}



\begin{quote}
Before the new one: say the Bengali for to ring, then the Bengali for a bearer.
\end{quote}

\subsection*{You'll want to know: \bn{মেনু}}
\textbf{\bn{মেনু}} — \emph{menu} — \textquotedblleft{}a menu\textquotedblright{}.

\begin{quote}
\bn{ওয়েটার} \bn{আমাদের} \bn{মেনু} \bn{দিলেন।}
\end{quote}

\begin{quote}
\emph{oyeṭār āmāder menu dilen.}
\end{quote}

\begin{quote}
The waiter gave us the menu.
\end{quote}

\begin{cousinweb}
\textbf{\bn{মেনু}} (\emph{menu}) is English \emph{menu}, itself from French.
\end{cousinweb}

\begin{grammarlens}[title={What you've built}]
One more word for getting around.
\end{grammarlens}

\subsection*{Your turn}
\begin{itemize}
\raggedright
  \item \emph{Say it:} \emph{menu}
  \item \emph{Say it:} \emph{menu}, once more
  \item \emph{Say it:} \emph{oyeṭār āmāder menu dilen.} — The waiter gave us the menu.
  \item \emph{Recall:} say the Bengali for to ring, then the Bengali for a bearer, then say \emph{menu} again
\end{itemize}

\begin{tcolorbox}[breakable,colback=teal!4,colframe=teal!35!black,title={Before you move on}]
What does \bn{মেনু} mean? (\textquotedblleft{}A menu\textquotedblright{}.) Say it once more.
\end{tcolorbox}
\section[\texorpdfstring{\bn{বকশিশ} (bôkshish)}{বকশিশ (bôkshish)}]{\bn{বকশিশ} (bôkshish) — a tip}

\label{lesson:BN-C407-bokshish}



\begin{quote}
Before the new one: say the Bengali for a bearer, then the Bengali for a menu.
\end{quote}

\subsection*{You'll want to know: \bn{বকশিশ}}
\textbf{\bn{বকশিশ}} — \emph{bôkshish} — \textquotedblleft{}a tip\textquotedblright{}.

\begin{quote}
\bn{আমরা} \bn{ওয়েটারকে} \bn{কুড়ি} \bn{টাকা} \bn{বকশিশ} \bn{দিলাম।}
\end{quote}

\begin{quote}
\emph{āmrā oyeṭārke kuṛi ṭākā bôkshish dilām.}
\end{quote}

\begin{quote}
We gave the waiter a twenty-taka tip.
\end{quote}

\begin{cousinweb}
\textbf{\bn{বকশিশ}} (\emph{bôkshish}) is usually traced to Persian \emph{bakhshish}, \textquotedblleft{}a gift, a present\textquotedblright{}, from \emph{bakhshīdan}, to give; English \emph{baksheesh} is the same word.
\end{cousinweb}

\begin{grammarlens}[title={What you've built}]
One more word for getting around.
\end{grammarlens}

\subsection*{Your turn}
\begin{itemize}
\raggedright
  \item \emph{Say it:} \emph{bôkshish}
  \item \emph{Say it:} \emph{bôkshish}, once more
  \item \emph{Say it:} \emph{āmrā oyeṭārke kuṛi ṭākā bôkshish dilām.} — We gave the waiter a twenty-taka tip.
  \item \emph{Recall:} say the Bengali for a bearer, then the Bengali for a menu, then say \emph{bôkshish} again
\end{itemize}

\begin{tcolorbox}[breakable,colback=teal!4,colframe=teal!35!black,title={Before you move on}]
What does \bn{বকশিশ} mean? (\textquotedblleft{}A tip\textquotedblright{}.) Say it once more.
\end{tcolorbox}
\section[\texorpdfstring{\bn{লন্ড্রি} (lônḍri)}{লন্ড্রি (lônḍri)}]{\bn{লন্ড্রি} (lônḍri) — a laundry}

\label{lesson:BN-C407-londri}



\begin{quote}
Before the new one: say the Bengali for a menu, then the Bengali for a tip (for service).
\end{quote}

\subsection*{You'll want to know: \bn{লন্ড্রি}}
\textbf{\bn{লন্ড্রি}} — \emph{lônḍri} — \textquotedblleft{}a laundry\textquotedblright{}.

\begin{quote}
\bn{হোটেলের} \bn{লন্ড্রি} \bn{আমার} \bn{জামা} \bn{কেচে} \bn{দিল।}
\end{quote}

\begin{quote}
\emph{hoṭeler lônḍri āmār jāmā keche dilô.}
\end{quote}

\begin{quote}
The hotel laundry washed my shirts.
\end{quote}

\begin{cousinweb}
\textbf{\bn{লন্ড্রি}} (\emph{lônḍri}) is borrowed from English \emph{laundry}; at home the same job is done by the \bn{ধোপা}.
\end{cousinweb}

\begin{grammarlens}[title={What you've built}]
One more word for getting around.
\end{grammarlens}

\subsection*{Your turn}
\begin{itemize}
\raggedright
  \item \emph{Say it:} \emph{lônḍri}
  \item \emph{Say it:} \emph{lônḍri}, once more
  \item \emph{Say it:} \emph{hoṭeler lônḍri āmār jāmā keche dilô.} — The hotel laundry washed my shirts.
  \item \emph{Recall:} say the Bengali for a menu, then the Bengali for a tip (for service), then say \emph{lônḍri} again
\end{itemize}

\begin{tcolorbox}[breakable,colback=teal!4,colframe=teal!35!black,title={Before you move on}]
What does \bn{লন্ড্রি} mean? (\textquotedblleft{}A laundry\textquotedblright{}.) Say it once more.
\end{tcolorbox}
\section[(dialogue)]{Looking back: from Addresses and Signs to Bells, Menus and Tips}

\label{lesson:BN-R407-review-400-407}



\begin{quote}
Say the Bengali for an address and a sign.
\end{quote}

\subsection*{Your turn}
Say these aloud:

\begin{itemize}
\raggedright
  \item an address, a sign, opposite, an instruction, to point (to)
  \item south, east, west, a compass, the horizon
  \item the middle, a detour, an edge, at a loss, to go round in circles
  \item a room (in a house), a booking, a lift, a reception desk, to arrange
  \item comfort, leisure, breakfast, to rest a while, to lean back
  \item a bathroom, a toilet, a water heater, an air conditioner, to clean
  \item a guesthouse, a hostel, a rest house for pilgrims, a place to stay, to tap (on a door)
  \item to ring, a bearer, a menu, a tip (for service), a laundry
\end{itemize}

\begin{tcolorbox}[breakable,colback=teal!4,colframe=teal!35!black,title={Before you move on}]
An address? (\textbf{\bn{ঠিকানা}}, \emph{ṭhikānā}.) A sign? (\textbf{\bn{চিহ্ন}}, \emph{chihnô}.) Opposite? (\textbf{\bn{উল্টো}}, \emph{ulṭo}.) An instruction? (\textbf{\bn{নির্দেশ}}, \emph{nirdesh}.) To point (to)? (\textbf{\bn{ইঙ্গিত} \bn{করা}}, \emph{inggit kôrā}.) South? (\textbf{\bn{দক্ষিণ}}, \emph{dôkkhin}.) East? (\textbf{\bn{পূর্ব}}, \emph{pūrbô}.) West? (\textbf{\bn{পশ্চিম}}, \emph{pôshchim}.) A compass? (\textbf{\bn{কম্পাস}}, \emph{kômpās}.) The horizon? (\textbf{\bn{দিগন্ত}}, \emph{digôntô}.) The middle? (\textbf{\bn{মাঝখান}}, \emph{mājhkhān}.) A detour? (\textbf{\bn{ঘুরপথ}}, \emph{ghurpôth}.) An edge? (\textbf{\bn{কিনারা}}, \emph{kinārā}.) At a loss? (\textbf{\bn{দিশেহারা}}, \emph{dishehārā}.) To go round in circles? (\textbf{\bn{ঘুরপাক} \bn{খাওয়া}}, \emph{ghurpāk khāoyā}.) A room (in a house)? (\textbf{\bn{ঘর}}, \emph{ghôr}.) A booking? (\textbf{\bn{বুকিং}}, \emph{buking}.) A lift? (\textbf{\bn{লিফট}}, \emph{liphṭ}.) A reception desk? (\textbf{\bn{রিসেপশন}}, \emph{risepshôn}.) To arrange? (\textbf{\bn{বন্দোবস্ত} \bn{করা}}, \emph{bôndobôstô kôrā}.) Comfort? (\textbf{\bn{আরাম}}, \emph{ārām}.) Leisure? (\textbf{\bn{আয়েশ}}, \emph{āyesh}.) Breakfast? (\textbf{\bn{জলখাবার}}, \emph{jôlkhābār}.) To rest a while? (\textbf{\bn{জিরোনো}}, \emph{jirono}.) To lean back? (\textbf{\bn{হেলান} \bn{দেওয়া}}, \emph{helān deoyā}.) A bathroom? (\textbf{\bn{বাথরুম}}, \emph{bāthrum}.) A toilet? (\textbf{\bn{টয়লেট}}, \emph{ṭôyleṭ}.) A water heater? (\textbf{\bn{গিজার}}, \emph{gijār}.) An air conditioner? (\textbf{\bn{এসি}}, \emph{esi}.) To clean? (\textbf{\bn{পরিষ্কার} \bn{করা}}, \emph{pôrishkār kôrā}.) A guesthouse? (\textbf{\bn{গেস্টহাউস}}, \emph{gesṭhāus}.) A hostel? (\textbf{\bn{হোস্টেল}}, \emph{hosṭel}.) A rest house for pilgrims? (\textbf{\bn{ধর্মশালা}}, \emph{dhôrmôshālā}.) A place to stay? (\textbf{\bn{ঠাঁই}}, \emph{ṭhā̃i}.) To tap (on a door)? (\textbf{\bn{টোকা} \bn{দেওয়া}}, \emph{ṭokā deoyā}.) To ring? (\textbf{\bn{বাজা}}, \emph{bājā}.) A bearer? (\textbf{\bn{বেয়ারা}}, \emph{beyārā}.) A menu? (\textbf{\bn{মেনু}}, \emph{menu}.) A tip (for service)? (\textbf{\bn{বকশিশ}}, \emph{bôkshish}.) A laundry? (\textbf{\bn{লন্ড্রি}}, \emph{lônḍri}.)
\end{tcolorbox}
